\section{WTO 法律制度}
\begin{enumerate}
  \item 1 个原则，2 个协议和 1 个争端解决机制。
  \item 以非歧视为目标，追求贸易的自由和公平。
  \item 基本知识。
        \begin{enumerate}
          \item 成员。主权国家和单独关税区。不能称成员国、缔约国，叫\zd{成员方、缔约方}。
          \item \zd{民用航空器贸易协议和政府采购协议}对中国没有约束力。
          \item 常设最高权力机构。总理事会。兼职\qd{争端解决机构和贸易政策审议}的职能。一套人马三个牌子。
          \item 我国义务 = 共同 + 特有。
        \end{enumerate}
  \item 基本原则。
        \begin{enumerate}
          \item \zd{国民待遇原则}。内外待遇一致，外国人待遇等于本国人。两方主体。仅限于列入承诺表的部门。
          \item \zd{最惠国待遇原则}。外国之间你最优惠。三方主体。
                \begin{enumerate}
                  \item 普遍性。给了一个给所有。
                  \item 自动性。多边无条件。允许搭便车。
                  \item 相互性。相互搭便车。
                  \item 同一性。同一状态延及不变形。
                  \item 例外。\yi 边贸优惠；\er 区域经济安排。\al{如欧盟}；\san 反倾销反补贴。
                  \item 任何协议的最惠国待遇条款，如果要修改需要全体成员同意。
                \end{enumerate}
        \end{enumerate}
  \item TRIMS 协议。与贸易有关的投资措施协议。
        \begin{enumerate}
          \item 4 种禁止性投资措施。\zd{当地成分要求、贸易平衡要求（进口<出口）、进口用汇限制、国内销售要求}。
          \item 核心。\zd{与货物贸易有关}。
        \end{enumerate}
  \item 服务贸易总协定。核心点在服务贸易中哪个要素跨境。
        \begin{kk}{分类}
          \begin{enumerate}
            \item 仅服务本身在跨境 \yjt 跨境交付/供应。
            \item 花钱的人跨境买服务 \yjt 境外消费。\al{中国某旅行社组团到泰国旅游}。
            \item 挣钱的人跨境提供服务 \yjt 挣钱存在。
            \item 设立机构挣钱 \yjt 商业存在。\al{加拿大某银行在中国设立分支机构}。
            \item 没有设立机构，自然人跑过去挣钱 \yjt \zd{自然人存在/流动}。\al{中国某运动员应聘到美国担任体育教练}。
          \end{enumerate}
        \end{kk}
        \begin{enumerate}
          \item 没有服务市场\zd{开放同一水平}的要求。
          \item \zd{依承诺}开放国际服务市场。
          \item \qd{开放服务市场，也必须满足最惠国待遇原则，不能厚此薄彼}。
        \end{enumerate}
  \item 争端的类型。
        \begin{enumerate}
          \item 违反性申诉。证明被诉方\zd{违反了哪个协议}。改 \hz 废。
          \item 非违反性申诉。证明被诉方的\zd{某一项法律法规 \hz 政策}导致利益受损 \hz 丧失。补偿。
          \item \zd{有的改就没的偿}。
        \end{enumerate}
  \item 程序。
        \begin{enumerate}
          \item 不以用尽当地救济为前提。
          \item 主体不包括企业。
          \item 磋商必经（仲裁、调解不是）。磋商前面加定语都是错的，\scx{充分磋商}。
                \begin{enumerate}
                  \item \zd{60} 天。过程、结果均\zd{保密}。
                  \item 仅是程序性要求，无实质影响（\scx{被用于后续争端解决}）。
                \end{enumerate}
          \item 专家组非常设，\zd{告啥审啥，不告不审}。
          \item 上诉机构常设，仅审查法律问题；无权发回重审；7 名成员中 3 人组成上诉庭。
          \item 争端解决机构通过\zd{报告}让其变成裁决。
          \item 投票采用\zd{反向一致}原则，一票赞成即通过。没有回避制度。
          \item 没有强制执行程序，裁决不履行，\zd{获权}方能报复。
          \item 允许交叉报复（不同种产品），但要\zd{与受损害的程度相等}。
          \item 必须做出裁决。
        \end{enumerate}
\end{enumerate}
